% !PRINTED 225
\weberpage{225}
%<<CONT
oder, wenn man mit $\omega'^2$ multiplicirt, nach (5) und (6)
\[
\begin{aligned}
\varepsilon^{r+s} + \varepsilon^{-(r+s)} - 1 &= -\varepsilon - \varepsilon^{-1} - 2 \\
\varepsilon + \varepsilon^{-1} + \varepsilon^{r+s} + \varepsilon^{-(r+s)} + 1 &= 0,
\end{aligned}
\]
eine Bedingung, die dann und nur dann erfüllt ist, wenn
$r + s \equiv \pm 2 \pmod{5}$ ist.

Wenn (14) von der dritten Ordnung sein soll, so muss
nach (18) und (5)
\[
\begin{aligned}
\varepsilon^2 + \varepsilon^{-2} + 2 &= 1 - \varepsilon^{-(r+s)} - \varepsilon^{r+s} \\
\varepsilon^2 + \varepsilon^{-2} + \varepsilon^{r+s} + \varepsilon^{-(r+s)} + 1 &= 0
\end{aligned}
\]
sein, und diese Bedingung ist dann und nur dann befriedigt,
wenn $r + s \equiv \pm 1 \pmod{5}$ ist.

Hiernach erhalten wir, wie es sein muss, in (13) und (14)
gerade 10 Substitutionen zweiter und 20 Substitutionen dritter
Ordnung, und die anderen Fälle bleiben also für die fünfte Ordnung
übrig.

Fassen wir das Resultat dieser Betrachtung zusammen, so
haben wir:

Die Substitution $\chi$ muss den Ausdruck haben:
\weq{19}{\chi = \begin{pmatrix} \omega, & 1 \\ 1, & -\omega \end{pmatrix},}
und unter den 60 Substitutionen (10):
\weq{20}{\Theta^r, \; \psi\Theta^r, \; \Theta^r\chi\Theta^s, \; \Theta^r\chi\psi\Theta^s}
kommen ausser der Identität vor:

15 \gesperrt{Substitutionen} $2^{\text{ten}}$ \gesperrt{Grades}:

$\psi\Theta^r, \; \Theta^r\chi\Theta^{-r}, \; \Theta^r\chi\psi\Theta^{-r}, \; r = 0, 1, 2, 3, 4,$
\weq{21}{
\begin{aligned}
&\text{20 }\text{\gesperrt{Substitutionen}}\ 3^{\text{ten}}\ \text{\gesperrt{Grades}:} \\
&\Theta^r\chi\Theta^s, \quad r+s \equiv \pm 2 \pmod{5} \\
&\Theta^r\chi\psi\Theta^s, \quad r+s \equiv \pm 1 \pmod{5}, \\
&\text{24 }\text{\gesperrt{Substitutionen}}\ 5^{\text{ten}}\ \text{\gesperrt{Grades}:} \\
&\Theta^r, \; \Theta^r\chi\Theta^s, \quad r+s \equiv \pm 1 \pmod{5} \\
&\Theta^r\chi\psi\Theta^s, \quad r+s \equiv \pm 2 \pmod{5}.
\end{aligned}
}
In expliciter Form erhält man für die Substitutionen (20)
den Ausdruck:
\weq{22}{
\begin{pmatrix} \varepsilon^r, & 0 \\ 0, & 1 \end{pmatrix}, \;
\begin{pmatrix} 0, & \varepsilon^r \\ -1, & 0 \end{pmatrix}, \;
\begin{pmatrix} \varepsilon^r\omega, & \varepsilon^{r-s} \\ 1, & -\varepsilon^{-s}\omega \end{pmatrix}, \;
\begin{pmatrix} -\varepsilon^{r+s}, & \varepsilon^r\omega \\ \varepsilon^s\omega, & 1 \end{pmatrix}
}
oder
\[
\varepsilon^r x, \quad \frac{-\varepsilon^r}{x}, \quad \frac{\varepsilon^r\omega x + \varepsilon^{r-s}}{x - \varepsilon^{-s}\omega}, \quad \frac{-\varepsilon^{r+s} x + \varepsilon^r\omega}{\varepsilon^s\omega x + 1}.
\]
